\section{Introduction}
A prepared experiment determines which future tests are executable, which histories have the same predictive content, and which distributions on that content are actually obtained. A finite observer must realize its predictions by one causal program. The geometry of a formal parameter space does not determine the cost of doing so. Neither does the rate at which an unrelated, cost-universal Markov chain mixes. The relevant obstruction is the portion of temporal degeneration that is visible to the task being performed.

This distinction becomes essential when physical preparation leaves the observer's memory intact. The preparation erases physical history but need not erase evidence about a fixed unknown parameter. Successive excursions are then conditionally Markov-renewal, rather than independent cycles. They may lead to several closed boundary classes with different durations. Their physical-time gains are ratios formed separately in each class and only then mixed by absorption probabilities. Near a recurrent-rank change, the entire ergodic projection can be discontinuous even though the prediction or decision under study remains uniformly well behaved.

The main result, Theorem~\ref{thm:transfer}, gives an exact criterion for this task-sensitive regularity. It does not postulate a spectral gap or a uniform bound for a group inverse. Under explicit compactness, finite-response continuity and return-uniform-integrability assumptions, it proves the equivalence of gain continuity, a compatibility condition at every limiting ergodic projection, uniform finite-acquisition convergence, uniform normalized-discount convergence, and a uniform approximate-corrector property. Thus it distinguishes a genuine average-risk discontinuity from a dynamically slow direction whose task amplitude vanishes. The criterion applies to all entering labels; a statement about one fortuitous initial distribution is weaker.

The quantitative part is based on the complete conditional excursion. If $P$ is its exit-label matrix, $T_{ij}=\E_i[\tau\one_{\{I_1=j\}}]$ its duration--exit matrix, $r$ its reward vector and $g$ its physical-time gain, the relevant centered reward is $r-Tg$. An approximate solution of
\begin{equation}\label{eq:intro-poisson}
 (I-P)u=r-Tg
\end{equation}
controls the discrepancy at a prescribed number of real physical calls. Duration and exit label may be correlated. This formulation avoids first paying a worst-case time to a recurrent root. Equation~\eqref{eq:intro-poisson} belongs to classical multichain semi-Markov bias theory; the present use is a uniform approximate-corrector criterion, its converse through recurrent-rank degeneration, and its coupling to one retained causal program and actual predictive occupation.

The result has three consequences relevant to the foundational problem. First, on an entire compact legal program class satisfying the criterion, a robust fixed-memory optimum is attained and is determined by finite model subfamilies. No recurrence constraint is inserted into that conclusion. Second, on finite algebraic instruments, task compatibility is decidable from the primitive model and yields an algebraic finite-acquisition modulus. Even without compatibility, a separate semialgebraic selection theorem supplies a near-optimal program whose recurrence cost grows at most polynomially in the target error, under a stated common-support hypothesis. Third, applying the criterion to bounded predictive tests transfers the actual occupation law. It can then be combined with a small-ball witness, a global candidate-valid cover and an actual-law update estimate to obtain matching causal and checkpoint resource curves.

The sharp examples are not substitutes for this theorem. They test it. In a rare-revelation decision experiment with revelation probability $|\theta|^k$, two retained labels are sufficient for vanishing worst-model risk, while one label is not. The exact finite-acquisition value is of order $N^{-1/k}$. The group inverse and the exact task bias are unbounded, but the approximate-corrector price has order $a^{-(k-1)}$, giving precisely this rate. A fixed noncommuting quantum instrument realizes the same task mechanism through a change of preparation rank. A different raw realization uses action-uniform drift for unbounded queue excursions. Neither realization is used in proving the general theorem.

For a single four-call scored protocol, the diagnostic task can be combined with a fresh geometric mark and a fixed inaccessible calibration sign. If the actually loaded mark has quantization error $q_M$, the optimized retained-memory value is bounded below and above by
\begin{equation}\label{eq:intro-law}
 q_M+\delta^2+D_n(e)
 \quad\hbox{and}\quad
 q_{\lfloor M/2\rfloor}+\delta^2+D_n(e),\qquad N=4n.
\end{equation}
Here $D_n(e)$ is an exact all-strategy diagnostic value and $e$ is a real reduction in revelation intensity, not an allowed tolerance silently promoted to a lower bound. Its common-task causal deficiency is comparable to $e$. Under acquired mass and cover bounds this gives
\[
 \calR(N,M,\delta,\varepsilon)
 \asymp M^{-2/d}+N^{-1/k}+\delta^2+\varepsilon^{1/k}.
\]
The fractional deficiency exponent is an attained task-sensitive amplification. General defect transfer, in contrast, requires regularity at both endpoints. The distinction is maintained throughout.

\subsection*{Relation to earlier theory}
The linear algebra of multichain evaluation is classical. Denardo and Fox~\cite{DenardoFox} treat multichain Markov-renewal programs; Federgrun and Schweitzer~\cite{FedergrunSchweitzer} study undiscounted Markov-renewal fixed points; Lamond and Puterman~\cite{LamondPuterman} develop generalized-inverse methods for Markov decision processes. Group inverses, classwise policy evaluation and Poisson equations are not new contributions here. In particular, replacing a correlated holding-time mark by its conditional mean independently of the exit state is not a permissible simplification of that theory.

The distinction in Theorem~\ref{thm:transfer} is between cost-universal recurrence and a specified continuous family of tasks. It gives a converse and a quantitative bridge on a compact family of raw causal responses, at an unchanged observer-state budget. Robust compact intersection is then a consequence, not a minimax interchange. Semialgebraic value asymptotics have important predecessors in Bewley and Kohlberg~\cite{BewleyKohlberg,BewleyKohlbergStationary}; real-algebraic choice and growth are standard~\cite{BasuPollackRoy}. We use them to certify task-compatible retained observers and to quantify a separate recurrence-exhaustion schedule. No efficient global synthesis algorithm is claimed.

Minimal predictive states and comparison of experiments have a long history~\cite{ShaliziCrutchfield,Blackwell,Fritz}. Our quotient lemma states its Borel realization assumptions rather than deriving them from positivity. Quantization and finite-memory learning are also established subjects~\cite{GrafLuschgy,HellmanCover}. The new resource statement identifies the actual occupation measure and the task-visible temporal price before applying those tools. The primitive instrument constructions below are deliberately explicit so that none of these layers can replace another by notation alone.
